\section[齊風]{齊風}

\subsection[雞鳴]{雞鳴}

雞既鳴矣，朝既盈矣。匪雞則鳴，蒼蠅之聲。

東方明矣，朝既昌矣。匪東方則明，月出之光。

蟲飛薨薨，甘與子同夢。會且歸矣，無庶予子憎！

\subsection[還]{還}

子之還兮，遭我乎峱之間兮。並驅從兩肩兮，揖我謂我儇兮。

子之茂兮，遭我乎峱之道兮。並驅從兩牡兮，揖我謂我好兮。

子之昌兮，遭我乎峱之陽兮。並驅從兩狼兮，揖我謂我臧兮。

\subsection[著]{著}

俟我於著乎而，充耳以素乎而，尚之以瓊華乎而。

俟我於庭乎而，充耳以青乎而，尚之以瓊瑩乎而。

俟我於堂乎而，充耳以黄乎而，尚之以瓊英乎而。

\subsection[東方之日]{東方之日}

東方之日兮，彼姝者子，在我室兮，在我室兮，履我即兮。

東方之月兮，彼姝者子，在我闥兮，在我闥兮，履我發兮。

\subsection[東方未明]{東方未明}

東方未明，顛倒衣裳。顛之倒之，自公召之。

東方未晞，顛倒裳衣。倒之顛之，自公令之。

折柳樊圃，狂夫瞿瞿。不能辰夜，不夙則莫。

\subsection[南山]{南山}

南山崔崔，雄狐綏綏。魯道有蕩，齊子由歸。既曰歸止，曷又懷止！

葛屨五兩，冠緌雙止。魯道有蕩，齊子庸止。既曰庸止，曷又從止！

蓺麻如之何？衡從其畝。取妻如之何？必告父母，既曰告止，曷又鞠止！

析薪如之何？匪斧不克。取妻如之何？匪媒不得，既曰得止，曷又極止！

\subsection[甫田]{甫田}

無田甫田，維莠驕驕。無思遠人，勞心忉忉。

無田甫田，維莠桀桀。無思遠人，勞心怛怛。

婉兮孌兮，緫角丱兮。未幾見兮，突而弁兮。

\subsection[盧令]{盧令}

盧令令，其人美且仁。

盧重環，其人美且鬈。

盧重鋂，其人美且偲。

\subsection[敝笱]{敝笱}

敝笱在梁，其魚魴鰥。齊子歸止，其從如雲。

敝笱在梁，其魚魴鱮。齊子歸止，其從如雨。

敝笱在梁，其魚唯唯。齊子歸止，其從如水。

\subsection[載驅]{載驅}

載驅薄薄，簟茀朱鞹。魯道有蕩，齊子發夕。

四驪濟濟，垂轡濔濔。魯道有蕩，齊子豈弟。

汶水湯湯，行人彭彭。魯道有蕩，齊子翱翔。

汶水滔滔，行人儦儦。魯道有蕩，齊子遊敖。

\subsection[猗嗟]{猗嗟}

猗嗟昌兮！頎而長兮。抑若揚兮，美目揚兮。巧趨蹌兮。射則臧兮。

猗嗟名兮！美目清兮。儀既成兮，終日射侯，不出正兮，展我甥兮。

猗嗟孌兮！清揚婉兮。舞則選兮，射則貫兮，四矢弁兮，以禦亂兮。
